\documentclass[11pt]{article}
\usepackage[margin=0.9in]{geometry}
\usepackage{amsmath,amssymb}
\usepackage{parskip}
\usepackage{fancyhdr}
\usepackage{enumitem}

\setlength{\headheight}{14pt}
\pagestyle{fancy}
\fancyhf{}
\lhead{CPSC 4710/5710 -- In-Class Work Session}
\rhead{Brainstorming Questions}
\cfoot{\thepage}

\begin{document}

\begin{center}
{\Large\bfseries Which Surrogates Matter?}\\[4pt]
{\large What Drives Success in Ensemble Transfer Attacks}
\end{center}

\textbf{Team members:} Mo Liu, Connor Kilrain, Tom Heim \hfill \textbf{Application area:} Computer Vision
\vspace{-6pt}

\hrulefill

\begin{enumerate}[leftmargin=*, itemsep=8pt]

\item \textbf{What task will you study? What dataset(s) and baseline model will you start with? Identify one primary trust axis and one secondary trust axis to address.}

Our team studies transfer-based black-box adversarial attacks on image classifiers, in which an attacker crafts adversarial examples on an ensemble of publicly available ``surrogate'' models and hopes they also fool an unseen target model; each of us takes a different angle on what makes some ensembles transfer better than others, and we will synthesize our findings. We will use standard image classification benchmarks (e.g., CIFAR-10, ImageNet) with a shared pool of publicly available pretrained classifiers, and a standard equal-weight ensemble attack as our common baseline. \textbf{Primary trust axis:} adversarial robustness; \textbf{secondary trust axis:} explainability.

\item \textbf{What is the proposed method you will implement, adapt, or change to address the problem?}

My part treats the surrogates as players in a cooperative game, where the payoff of any subset is how often an ensemble attack built from it fools held-out targets; each surrogate's Shapley value is then its average contribution, and pairwise interaction values show which surrogates are redundant and which are complementary. Since the number of subsets grows exponentially, I will compute these values exactly for a small pilot pool and estimate them by sampling for larger ones. I will then test whether explanation similarity between a surrogate and the target (e.g., how alike their saliency maps are) predicts Shapley value better than simple factors like architecture type or accuracy, and use the best predictor to pick a small, effective set of surrogates without rerunning the full analysis.

\item \textbf{What evaluation metric will you use, and what stress test or sanity check will you run?}

My main metric is the attack success rate on held-out targets, along with how well each predictor tracks the surrogates' Shapley values. To assess efficiency, I will measure how much of the full ensemble's success rate a small selected subset retains at a fraction of the compute; this builds on Shapley's efficiency axiom, under which the full ensemble's success (above the clean-image baseline) is fully distributed across the surrogates, so each Shapley value reads as that surrogate's share of the credit. As a stress test, I will evaluate on targets whose architecture type is absent from the surrogate pool; as a sanity check, I will confirm that a deliberately useless surrogate (e.g., an untrained model) receives near-zero credit, as the Shapley axioms require.

\item \textbf{What are the main data, compute, or timeline constraints? What is the biggest risk and your fallback plan?}

The main constraint is compute, since every part of the project needs many attack runs, so we will share one attack pipeline and model pool and pilot on a small dataset before scaling up. The biggest risk is that one simple factor, such as architecture type, explains most differences in transfer; our fallback is to report how the factors compare, which still shows which surrogate properties matter most. Each part also stands on its own if another does not pan out.

\end{enumerate}

\end{document}
